% segment-local-id: OR03-S009
% segment-id: d90.orig.v1.tr03.seg0009
% license: CC BY-SA 4.0
% provenance: Materi asesmen asli berbahasa Indonesia, disusun secara independen.

\section{Rubrik analitik untuk pembuktian}
\label{orig03:rubric}

Setiap kriteria di bawah dinilai dengan skala yang sama:
\textbf{0} tidak ada atau tidak sah;
\textbf{1} fragmen relevan dengan galat mayor;
\textbf{2} arah sebagian benar tetapi masih ada celah substantif;
\textbf{3} benar dengan satu kekurangan minor yang tidak mengubah kesimpulan;
dan \textbf{4} lengkap, tepat, dan seluruh dependensi logis dinyatakan.
Setiap rubrik mempunyai empat kriteria, jadi skor maksimumnya adalah 16.
Keanggunan gaya tidak boleh menutupi galat matematika, dan notasi alternatif
yang ekuivalen harus dinilai berdasarkan isinya.

% stable-id: d90.orig.v1.tr03.rubric.proof.0001
\subsection{Rubrik proof.0001: hipotesis, tipe, dan kuantor}
\label{orig03:rubric:proof:0001}

\paragraph{Kriteria (masing-masing 0--4).}
\begin{enumerate}
  \item Menyatakan objek, ruang, domain, dan kodomain secara konsisten.
  \item Menempatkan semua kuantor dan dependensi parameter pada lingkup yang
  benar, termasuk perbedaan antara pernyataan titik demi titik dan seragam.
  \item Menggunakan tepat hipotesis ketertutupan, konveksitas, keterbatasan,
  atau keterukuran pada langkah yang memerlukannya.
  \item Menutup semua kasus tepi dan memastikan bahwa kesimpulan mempunyai
  tipe logis yang sama dengan klaim semula.
\end{enumerate}

\paragraph{Kegagalan umum.}
Menganggap domain sama dengan seluruh ruang; menukar ``untuk setiap'' dengan
``terdapat''; memakai konstanta yang diam-diam bergantung pada iterasi; membagi
dengan besaran yang belum dibuktikan positif; atau menyatakan limit tanpa
menentukan topologi.

% stable-id: d90.orig.v1.tr03.rubric.proof.0002
\subsection{Rubrik proof.0002: keberadaan dan argumen topologis}
\label{orig03:rubric:proof:0002}

\paragraph{Kriteria (masing-masing 0--4).}
\begin{enumerate}
  \item Mengidentifikasi himpunan sublevel atau barisan peminimum yang tepat
  dan membuktikan keterbatasan atau kekompakannya.
  \item Menggunakan semikontinuitas-bawah, ketertutupan, atau kekompakan dengan
  arah ketaksamaan dan topologi yang benar.
  \item Membedakan dengan tegas keberadaan, ketunggalan, dan kestabilan, serta
  memberi hipotesis tersendiri bagi masing-masing klaim.
  \item Menangani fungsi bernilai-diperluas, domain efektif, dan nilai tak
  hingga tanpa operasi aljabar yang tidak sah.
\end{enumerate}

\paragraph{Kegagalan umum.}
Mengatakan ``kontinu maka minimum ada'' pada himpunan tak kompak; menganggap
konveksitas memberi ketunggalan; menyimpulkan koersivitas hanya dari
keterbatasan-bawah; mengabaikan domain kosong; atau mengambil limit dalam
domain yang belum dibuktikan tertutup.

% stable-id: d90.orig.v1.tr03.rubric.proof.0003
\subsection{Rubrik proof.0003: penurunan, teleskop, dan laju}
\label{orig03:rubric:proof:0003}

\paragraph{Kriteria (masing-masing 0--4).}
\begin{enumerate}
  \item Menurunkan ketaksamaan satu-langkah dari asumsi yang dinyatakan, dengan
  tanda, faktor, dan norma yang benar.
  \item Menentukan besaran potensial nonnegatif dan meneleskopkan indeks tanpa
  kehilangan suku batas.
  \item Memilih langkah atau parameter dalam rentang yang benar dan melacak
  semua konstanta hingga batas akhir.
  \item Membedakan laju nilai fungsi, jarak iterasi, dan residu, termasuk apakah
  klaim berlaku untuk iterasi terakhir, terbaik, atau rata-rata.
\end{enumerate}

\paragraph{Kegagalan umum.}
Menghapus suku bertanda tak diketahui; kesalahan indeks awal/akhir; memilih
langkah setelah melihat kuantitas acak masa depan; mengubah batas
$O(k^{-1/2})$ menjadi $O(k^{-1})$ dengan menguadratkan sepihak; atau
menyembunyikan konstanta yang bergantung pada dimensi.

% stable-id: d90.orig.v1.tr03.rubric.proof.0004
\subsection{Rubrik proof.0004: dualitas, kualifikasi, dan KKT}
\label{orig03:rubric:proof:0004}

\paragraph{Kriteria (masing-masing 0--4).}
\begin{enumerate}
  \item Membentuk Lagrangian dengan tanda, ruang dual, dan domain yang benar.
  \item Menyatakan kondisi kualifikasi yang benar-benar cukup sebelum
  menggunakan dualitas kuat atau keberadaan pengali.
  \item Menurunkan kelayakan primal, kelayakan dual, stasioneritas, dan
  kelonggaran komplementer tanpa melewatkan inklusi subdiferensial.
  \item Membuktikan kebutuhan atau kecukupan kondisi KKT sesuai kelas masalah,
  serta menghubungkannya dengan titik pelana atau nilai dual.
\end{enumerate}

\paragraph{Kegagalan umum.}
Menyatakan KKT tanpa kualifikasi; memberi pengali bertanda salah; mengganti
inklusi dengan persamaan gradien saat fungsi tak mulus; menyamakan dualitas
lemah dengan kuat; atau memakai kelonggaran komplementer untuk menyimpulkan
kelayakan yang belum dibuktikan.

% stable-id: d90.orig.v1.tr03.rubric.proof.0005
\subsection{Rubrik proof.0005: filtrasi dan argumen stokastik}
\label{orig03:rubric:proof:0005}

\paragraph{Kriteria (masing-masing 0--4).}
\begin{enumerate}
  \item Mendefinisikan filtrasi dan menyatakan dengan benar besaran yang
  terukur sebelum pengambilan sampel berikutnya.
  \item Menghitung ekspektasi bersyarat, ketakbiasan, dan bias tanpa
  mengasumsikan independensi yang tidak tersedia.
  \item Membedakan batas varians dari batas momen kedua dan menggunakan hukum
  varians total atau Jensen secara sah.
  \item Menukar ekspektasi, penjumlahan, limit, atau subdiferensial hanya jika
  syarat integrabilitas dan teorema yang dibutuhkan telah dipenuhi.
\end{enumerate}

\paragraph{Kegagalan umum.}
Mengondisikan pada informasi yang sudah memuat sampel baru; menyebut galat
sebagai selisih martingal padahal mean bersyaratnya tidak nol; mengganti
$\mathbb E\lVert G\rVert^2$ dengan varians; memakai sampel tanpa pengembalian
seolah-olah independen; atau mengambil ekspektasi dari kuantitas tak
terintegralkan.

% stable-id: d90.orig.v1.tr03.rubric.proof.0006
\subsection{Rubrik proof.0006: operator monoton dan resolven}
\label{orig03:rubric:proof:0006}

\paragraph{Kriteria (masing-masing 0--4).}
\begin{enumerate}
  \item Menyatakan domain operator bernilai-himpunan dan membedakan monotoni,
  monotoni maksimal, monotoni kuat, serta kokersivitas.
  \item Mengubah inklusi menjadi persamaan resolven atau titik tetap dengan
  urutan komposisi dan parameter yang benar.
  \item Menggunakan firm non-ekspansivitas, keaveragan, atau kontraksi dengan
  konstanta yang dibuktikan, bukan diasumsikan.
  \item Menghubungkan residu atau titik klaster dengan nol operator melalui
  ketertutupan graf dan argumen konvergensi yang lengkap.
\end{enumerate}

\paragraph{Kegagalan umum.}
Menganggap semua operator monoton maksimal; memperlakukan resolven sebagai
invers biasa; menukar $R_AR_B$ dengan $R_BR_A$; mengabaikan domain pada aturan
jumlah; atau menyimpulkan konvergensi iterasi hanya dari residu yang menuju nol
tanpa kontrol titik klaster.

% stable-id: d90.orig.v1.tr03.rubric.proof.0007
\subsection{Rubrik proof.0007: contoh penyangkal dan sertifikat komputasional}
\label{orig03:rubric:proof:0007}

\paragraph{Kriteria (masing-masing 0--4).}
\begin{enumerate}
  \item Menargetkan tepat satu hipotesis atau kesimpulan dan memastikan contoh
  memenuhi semua asumsi lain yang dipertahankan.
  \item Menghitung objek kunci secara eksak atau dengan batas galat yang
  tersertifikasi, termasuk kasus tepi yang relevan.
  \item Memisahkan bukti matematis dari ilustrasi numerik serta menyatakan
  presisi, toleransi, dan residu bila komputasi dipakai.
  \item Menjelaskan secara eksplisit bagaimana hasil contoh membatalkan klaim
  universal atau bagaimana sertifikat memverifikasi klaim yang terbatas.
\end{enumerate}

\paragraph{Kegagalan umum.}
Contoh melanggar asumsi lain; satu plot diperlakukan sebagai bukti universal;
residu kecil disamakan dengan galat solusi tanpa batas kondisi; pembulatan
mengubah tanda; atau program menguji hanya satu parameter sementara klaimnya
untuk semua parameter.

